\section{The half-open layer partition}
\label{sec:area_law_layer_partition}

The nested neighborhoods $N_k$ determine the layers
$D_k=N_{k+1}\setminus N_k$. Their half-open convention gives a
unique layer assignment before taking closures of birth regions;
see \cite[Section~11]{OpenAI2026AreaLaw}.

% Source: 10-geometry.tex, lines 154--181. All disjointness statements
% concern the half-open sets, not their closures.

\begin{theorem}[An actual fine cell belongs to its layer]
    \label{thm:area_law_fine_cell_containment}
    \lean{TNLean.PEPS.AreaLaw.Geometry.dyadicCell_subset_dyadicLayer_of_mem_fineLayerIndices}
    \leanok
    \uses{def:area_law_fine_layer_indices}
    For arbitrary origin, finite endpoint set and nonnegative integer
    width, suppose that $\ell\le k$ and that $z$ is an actual fine-layer
    index at exponent $\ell$. Then its half-open cell belongs to $D_k$.
\end{theorem}
\begin{proof}\leanok
    \uses{thm:area_law_fine_layer_count}
    The exact union of the fine cells in $D_k$ contains the cell
    indexed by $z$.
\end{proof}

\begin{theorem}[Pairwise disjoint half-open layers]
    \label{thm:area_law_layer_pairwise_disjoint}
    \lean{TNLean.PEPS.AreaLaw.Geometry.dyadicLayer_pairwiseDisjoint}
    \leanok
    \uses{def:area_law_dyadic_layers}
    For arbitrary origin, finite endpoint set and nonnegative integer
    width, $D_k\cap D_h=\varnothing$ whenever $k\ne h$.
\end{theorem}
\begin{proof}\leanok
    \uses{def:area_law_dyadic_layers,thm:area_law_dyadic_nesting}
    If $k<h$, then $D_k\subseteq N_{k+1}\subseteq N_h$, whereas
    $D_h$ is disjoint from $N_h$. The other order follows by symmetry.
\end{proof}

\begin{theorem}[The dummy region is disjoint from later layers]
    \label{thm:area_law_dummy_layer_disjoint}
    \lean{TNLean.PEPS.AreaLaw.Geometry.dyadicNeighborhood_disjoint_later_layer}
    \leanok
    \uses{def:area_law_dyadic_layers}
    Under the same hypotheses, fix a lower index $k_0$. For every
    $k\ge k_0$, the dummy neighborhood $N_{k_0}$ is disjoint from $D_k$.
\end{theorem}
\begin{proof}\leanok
    \uses{def:area_law_dyadic_layers,thm:area_law_dyadic_nesting}
    Nesting gives $N_{k_0}\subseteq N_k$, and
    $D_k=N_{k+1}\setminus N_k$.
\end{proof}

\begin{theorem}[Unique layer outside the dummy region]
    \label{thm:area_law_unique_layer}
    \lean{TNLean.PEPS.AreaLaw.Geometry.exists_unique_layer_of_not_mem_dyadicNeighborhood}
    \leanok
    \uses{def:area_law_dyadic_layers}
    Suppose that $Z$ is nonempty and $C_0\ge2$. For every lower
    index $k_0$ and every $x\notin N_{k_0}$, there is a unique
    $k\ge k_0$ such that $x\in D_k$.
\end{theorem}
\begin{proof}\leanok
    \uses{thm:area_law_dyadic_exhaustion,thm:area_law_dyadic_nesting,
        thm:area_law_layer_pairwise_disjoint}
    Exhaustion gives a least $K$ for which $x\in N_K$.
    Nesting and $x\notin N_{k_0}$ imply $K>k_0$. Minimality gives
    $x\notin N_{K-1}$, so $x\in D_{K-1}$. Pairwise disjointness
    gives uniqueness.
\end{proof}

\begin{theorem}[Exact cover by the dummy region and later layers]
    \label{thm:area_law_dummy_layers_cover}
    \lean{TNLean.PEPS.AreaLaw.Geometry.dyadicNeighborhood_union_iUnion_layers_eq_univ}
    \leanok
    \uses{def:area_law_dyadic_layers}
    For nonempty $Z$, integer $C_0\ge2$ and every lower index $k_0$,
    \begin{align}
        \mathbb R^2=N_{k_0}\cup\bigcup_{k\ge k_0}D_k.
    \end{align}
\end{theorem}
\begin{proof}\leanok
    \uses{thm:area_law_unique_layer}
    A point either belongs to $N_{k_0}$ or lies in its unique later
    layer.
\end{proof}
